This thesis asks how representation-aware execution and workload-level
refinement calibration can make JHQ an efficient, memory-resident GPU ANN
system. Representation-aware execution exploits the code structure in
data layout and scoring; workload-level calibration chooses one reusable
refinement budget from unlabelled queries. Construction and capacity
establish the system's foundation, while cross-system comparisons assess
its practical value. Four research questions follow.

\paragraph{RQ1: GPU index construction and capacity.}
Can a complete JHQ index be built and kept resident on one GPU, and at
what construction time, peak memory and resident capacity (\cref{sec:7-7})?

\paragraph{RQ2: Representation-aware execution.}
How can JHQ's code structure guide GPU table lookup, memory access and
candidate selection to improve throughput while maintaining each stage's
specified ranking and tie rules (\cref{sec:7-3,sec:7-4,sec:7-6})?

\paragraph{RQ3: Workload-level refinement calibration.}
Can a refinement budget calibrated once per workload without ground-truth
labels keep recall on separate test queries close to that of a larger-budget
search and repay its calibration cost (\cref{sec:7-5})?

\paragraph{RQ4: Cross-system search performance.}
Under host-to-host timing, at which recall targets and batch sizes does
JHQ-GPU outperform the tested GPU ANN baselines (\cref{sec:7-2})?

Chapter 3 turns these questions into controlled comparisons. Together they
separate what was built, whether its execution is correct, what effort a
workload requires, and whether the resulting system has useful operating
regions.
